\documentclass[11pt]{article}

\usepackage[margin=1in]{geometry}
\usepackage{amsmath,amssymb,amsthm}
\usepackage{booktabs}
\usepackage{hyperref}
\usepackage{enumitem}

\newtheorem{definition}{Definition}
\newtheorem{proposition}{Proposition}
\newtheorem{remark}{Remark}

\title{Representation Genesis: Wisdom-Oriented Agents Invent Better Languages for Future Thought}
\author{Mian Zhang}
\date{}

\begin{document}
\maketitle

\begin{abstract}
AI systems are usually optimized to produce better answers, policies, or single-shot benchmark scores.
We study a complementary objective: producing better representations for future thought.
We call this problem \emph{Representation Genesis}.
In this framework, macros, ontologies, theorem templates, task schemas, skill APIs, and evidence gates are treated as evolvable objects.
A representation is valuable only if it compresses future problems after paying its unpacking cost, transfers across tasks, remains robust under perturbation, and preserves provenance.
We introduce a representation score combining compression gain, unpacking cost, reuse, transfer, robustness, novelty, and residual absorption.
We report the first reproducible artifacts: a toy monoid macro-search lab, a 130-file Wisdom Science macro-mining pass, a 10-claim machine-readable registry, and generated LaTeX tables used by this draft.
We further introduce \emph{Perspectival Grounding}: the view is not the world, so a wise representation must preserve competing interpretations and select evidence-gathering views that collapse ambiguity without prematurely collapsing truth.
\end{abstract}

\section{Introduction}

Current AI evaluation asks whether a model can answer the present problem.
Wisdom Science asks a different question: after experience, failure, and feedback, does the system become better equipped for future problems?
Representation Genesis is the upstream version of this question.
Instead of only changing answers, a wise system changes the language in which future answers become shorter, safer, and more transferable.

\paragraph{Thesis.}
Wisdom-oriented agents should search over representations, not only over outputs.
The relevant product of a learning system is sometimes a macro, task schema, proof route, failure ontology, or API boundary that makes many future tasks easier.

\paragraph{Boundary.}
This paper does not claim that compression, macro discovery, or library learning are new.
The contribution is narrower: a governance framework for representation search inside a wisdom-oriented agent, with explicit evidence gates and transfer-oriented metrics.

\section{Motivation}

The Wisdom Science corpus already contains repeated conceptual structure: WisdomBench, Cognitive Immunity, Evidence Gates, the Intelligence-Wisdom Gap, and the Second Scaling Law.
This repetition is not merely stylistic.
It means that the research program has developed a reusable language.
If that language is left implicit, papers drift, formulas fork, and claims become hard to audit.
If it is made explicit, the same language becomes a compression layer for future research.

\section{Problem Formulation}

\begin{definition}[Representation Object]
A representation object $r$ is any reusable structure that changes how future tasks are encoded, solved, or audited.
Examples include macros, ontologies, theorem templates, task schemas, skill APIs, prompt schemas, failure taxonomies, and evidence gates.
\end{definition}

\begin{definition}[Representation Genesis]
Representation Genesis is the process of proposing, evaluating, retaining, and revising representation objects under compression, transfer, robustness, novelty, and provenance constraints.
\end{definition}

\begin{definition}[Representation Score]
For a representation $r$ evaluated on corpus $\mathcal{C}$ and held-out tasks $\mathcal{H}$, define
\[
S(r)=
\lambda_c G_c(r)
-\lambda_u U(r)
+\lambda_\rho R(r)
+\lambda_\tau T(r)
+\lambda_h H(r)
+\lambda_n N(r)
+\lambda_a A(r),
\]
where $G_c$ is compression gain, $U$ is unpacking cost, $R$ is reuse, $T$ is transfer gain, $H$ is robustness, $N$ is novelty distance, and $A$ is residual absorption.
\end{definition}

\paragraph{Compression gain.}
A representation must shorten a corpus after paying definition cost:
\[
G_c(r)=L(\mathcal{C})-[L(r)+L(\mathcal{C}\mid r)].
\]

\paragraph{Unpacking cost.}
High compression is not useful if the representation becomes too expensive to explain or verify.

\paragraph{Transfer gain.}
A representation should improve held-out tasks, not merely memorize the source corpus.

\paragraph{Residual absorption.}
Failures are especially valuable when they expose parts of the problem space that the current representation cannot compress.

\section{Algorithmic Loop}

Representation Genesis follows a simple loop:
\begin{enumerate}[leftmargin=*]
  \item \textbf{Propose}: generate a candidate macro, schema, ontology, or evidence rule.
  \item \textbf{Pack}: measure whether it compresses a source corpus after paying definition cost.
  \item \textbf{Unpack}: test whether another system can recover the intended meaning.
  \item \textbf{Transfer}: evaluate held-out tasks, papers, or code modules.
  \item \textbf{Stress}: perturb wording, task distribution, or evidence source.
  \item \textbf{Retain}: keep a Pareto frontier rather than a single winner.
\end{enumerate}

\begin{proposition}[Why a Pareto frontier is necessary]
If representations are scored only by raw compression, then the optimum may hide assumptions or destroy transfer.
Therefore a wisdom-oriented representation archive must preserve non-dominated candidates across compression, unpacking cost, transfer, robustness, novelty, and provenance.
\end{proposition}

\section{Reproducible Artifacts}

We implement the first version as dependency-free scripts:
\begin{itemize}[leftmargin=*]
  \item \texttt{experiments/representation\_lab\_toy\_monoid.py};
  \item \texttt{experiments/representation\_lab\_metrics.py};
  \item \texttt{experiments/wisdom\_science\_macro\_mining.py};
  \item \texttt{experiments/wisdom\_science\_claim\_registry.py};
  \item \texttt{experiments/generate\_wisdom\_science\_tables.py}.
\end{itemize}

\paragraph{Toy monoid.}
The toy lab verifies that macros compress only real source substrings.
It explicitly avoids the failure mode in which generated macro symbols are recursively compressed as if they were new source structure.
The current demo compresses a 53-symbol corpus to 48 symbols after definition cost, yielding a compression ratio of $1.10$ with one real macro.

\paragraph{Corpus macro mining.}
The macro-mining pass scans 130 paper and code files.
The strongest seed terms are shown in Table~\ref{tab:seed_terms}.

\begin{table}[t]
\centering
\small
\caption{Top seed terms from the Wisdom Science corpus. Counts and document frequencies are generated from \texttt{macro\_registry\_v0.json}.}
\label{tab:seed_terms}
\begin{tabular}{lrr}
\toprule
Term & Count & Doc. freq. \\
\midrule
immunity & 502 & 33 \\
WisdomBench & 178 & 24 \\
Cognitive Immunity & 144 & 27 \\
Ouroboros & 106 & 31 \\
homeostasis & 75 & 19 \\
plasticity & 64 & 14 \\
Second Scaling Law & 52 & 13 \\
Wisdom Quotient & 49 & 22 \\
Cognitive Entropy & 40 & 9 \\
Intelligence-Wisdom Gap & 35 & 12 \\
Evidence Gate & 32 & 13 \\
WisdomBench-Embodied & 26 & 10 \\
\bottomrule
\end{tabular}

\end{table}

\paragraph{Representation-compression study.}
We also run a deterministic macro-rewriting experiment on 12 core Wisdom Science source files.
The canonical macro glossary contains 13 long-form terms such as Wisdom Quotient, Second Scaling Law, Cognitive Immunity, WisdomBench-Embodied, and Representation Genesis.
After paying the glossary cost, the macro-coded corpus saves 3079 bytes and yields a compression ratio of $1.0126$.
The result is intentionally modest: it is evidence that a shared representation language has measurable compression value, not evidence that compression alone is wisdom.

\begin{table}[t]
\centering
\small
\caption{Canonical macro counts in the representation-compression study. Generated from \texttt{representation\_compression\_v0.json}.}
\label{tab:representation_compression}
\begin{tabular}{llr}
\toprule
Macro & Long form & Count \\
\midrule
WB & WisdomBench & 50 \\
CI & Cognitive Immunity & 30 \\
SSL & Second Scaling Law & 27 \\
EG & Evidence Gate & 25 \\
RG & Representation Genesis & 20 \\
IWG & Intelligence-Wisdom Gap & 18 \\
WB-E & WisdomBench-Embodied & 17 \\
WQ & Wisdom Quotient & 16 \\
EFI & Embodied Failure Immunity & 11 \\
COS & Cognitive Operating System & 5 \\
\bottomrule
\end{tabular}

\end{table}

\paragraph{Claim registry.}
We export the core claims into JSON, CSV, Markdown, and LaTeX tables.
This prevents a common failure mode in ambitious research programs: claims become stronger as they move from notes to papers.
The registry records claim, evidence source, strength, and limitation.

\begin{table}[t]
\centering
\scriptsize
\caption{Machine-generated claim registry used to keep Wisdom Science claims auditable.}
\label{tab:claim_registry}
\begin{tabular}{p{0.08\linewidth}p{0.46\linewidth}p{0.18\linewidth}p{0.14\linewidth}}
\toprule
ID & Claim & Evidence & Strength \\
\midrule
C1 & Wisdom should be distinguished from first-round intelligence. & P2, P3 & strong \\
C2 & Wisdom can be operationalized as normalized cross-round learning gain. & P1, P2 & strong \\
C3 & Higher first-round intelligence does not guarantee higher wisdom. & P3, P4 & medium-strong \\
C4 & Wisdom scales with architectural qualities and experience, not only parameters. & P4 & medium \\
C5 & Failures can be converted into transferable immunity signals. & P1 & medium-strong \\
C6 & Embodied agents should report longitudinal learning ability. & P5, P6, P7 & strong \\
C7 & Robot learning ability can be described by embodied plasticity, immunity, and homeostasis. & P5, P6 & medium \\
C8 & Evidence gates are necessary for embodied wisdom claims. & P7 & strong \\
C9 & Representation search is the next layer of wisdom-oriented systems. & P8 draft & new \\
C10 & Embodied failure trajectories can become transferable antigens. & P9 draft & new \\
C11 & A single observation should be treated as a projection rather than the latent world. & P8 toy lab & new \\
C12 & Embodied wisdom is distributed across organ-like subsystems, not only a brain policy. & P9 schema & new \\
\bottomrule
\end{tabular}

\end{table}

\section{Perspectival Grounding}

Representation Genesis is not only about shorter language.
It is also about choosing a language that does not confuse a projection with the world.
A glyph that appears to be ``16'' may become ``91'' under rotation or may belong to another writing system.
A Chinese character may carry multiple senses, emotions, and social uses until collocation and discourse context are observed.
An occluded object may have several affordances, and a strategic signal may be literal, a bluff, or a second-order feint.

\begin{definition}[Perspectival Grounding]
Let $x$ be a latent state, $v$ a viewpoint or context, and $o=\pi_v(x)$ an observation.
Perspectival Grounding is the discipline of preserving the hypothesis set
\[
\mathcal{H}(o)=\{x':\pi_v(x')=o\}
\]
until additional views, actions, context, or social evidence reduce ambiguity.
\end{definition}

This turns ambiguity from a nuisance into a representation objective.
The agent should not merely answer from the first projection; it should ask which next view, touch, query, or belief-model update has the best information gain per cost.

\begin{table}[t]
\centering
\small
\caption{Perspectival Grounding toy lab. Each case starts from a single ambiguous projection and chooses low-cost additional evidence until the latent hypothesis set is resolved. Generated from \texttt{perspectival\_grounding\_v0.json}.}
\label{tab:perspectival_grounding}
\begin{tabular}{p{0.25\linewidth}p{0.15\linewidth}rrrr}
\toprule
Case & Domain & $|\mathcal{H}_0|$ & $H_0$ & Drop & Cost \\
\midrule
rotated 16 91 & glyph & 3 & 1.58 & 1.58 & 1.0 \\
six nine orientation & glyph & 2 & 1.00 & 1.00 & 0.5 \\
polysemous chinese xing & language & 4 & 2.00 & 2.00 & 0.5 \\
occluded mug can bowl & robotics & 3 & 1.58 & 1.58 & 1.5 \\
deictic instruction & social\_robotics & 3 & 1.58 & 1.58 & 0.8 \\
counter reflexive game & social\_game & 3 & 1.58 & 1.58 & 1.2 \\
\midrule
Mean & -- & -- & 1.56 & 1.56 & -- \\
\bottomrule
\end{tabular}

\end{table}

The current dependency-free lab contains six cases across glyphs, language, robotics, deictic social instruction, and counter-reflexive game signals.
The mean initial ambiguity is 1.56 bits, the mean final ambiguity is 0 bits, and every case is resolved after evidence-gathering views are added.
This is not a claim that the toy lab solves perception.
It is a guardrail: any future wisdom representation must distinguish surface observation, latent state, viewpoint, social context, and evidence cost.

\section{Evidence Path}

Representation Genesis is designed as a toy-to-real evidence chain:
\begin{enumerate}[leftmargin=*]
  \item \textbf{Toy monoid lab}: verify that proposed macros compress after definition cost and do not recursively compress generated symbols.
  \item \textbf{Research corpus}: mine the Zenodo/P0--P7 corpus for repeated claims, formulas, evidence gates, and terminology drift.
  \item \textbf{Code corpus}: mine SOVEREIGN modules for reusable skill and pipeline abstractions.
  \item \textbf{Benchmark schemas}: rewrite WisdomBench and WB-E task schemas using the discovered representation language, then measure prompt cost and consistency.
  \item \textbf{Perspectival grounding}: test whether ambiguous observations retain multiple hypotheses until active views, context, or counter-reflexive belief models resolve them.
\end{enumerate}

\section{Relation to Wisdom Science}

Representation Genesis sits above the existing Wisdom Science stack.
WisdomBench measures whether agents improve across repeated experience.
Cognitive Immunity studies whether failures become transferable rules.
The Second Scaling Law studies how wisdom depends on architecture.
Representation Genesis asks a more upstream question: can a system invent the language that makes all of these future measurements easier to express, verify, and extend?

\section{Related Work}

\paragraph{Compression and mathematics.}
Freedman et al. frame mathematical abstraction as compression over formal structure.
Our work uses this as motivation but does not claim priority over compression accounts of mathematics.

\paragraph{Library learning and program synthesis.}
DreamCoder, Stitch, and LILO learn reusable program abstractions.
Representation Genesis differs by treating papers, benchmarks, evidence gates, failure taxonomies, and skill APIs as representation objects, not only programs.

\paragraph{Compression and intelligence.}
Schmidhuber and Hutter connect compression, curiosity, and universal intelligence.
Our claim is not that compression is intelligence, but that representation governance is one measurable component of wisdom-oriented systems.

\paragraph{Symbolic law discovery.}
AI Feynman and symbolic-regression systems search compact laws.
Representation Genesis searches a broader class of research language: formulas, macros, schemas, and provenance rules.

\paragraph{Active perception and symbol grounding.}
Active perception and next-best-view robotics study how agents gather information under partial observability, while symbol grounding studies how signs acquire meaning through perception and action.
Perspectival Grounding does not claim priority over those fields.
It contributes a Wisdom Science formulation: ambiguity, viewpoint, semantic context, and counter-reflexive belief should be reported as evidence objects rather than hidden inside a single model answer.

\section{Limitations}

Compression is not truth.
A shorter representation may hide assumptions, overfit to a corpus, or reduce interpretability.
The current macro miner is lexical rather than semantic.
The current toy monoid is a sanity check, not evidence about real mathematics.
Future work must test whether representation changes improve held-out writing, benchmark design, and embodied failure transfer.
The Perspectival Grounding lab is intentionally toy-scale; real use requires multimodal perception logs, active sensing traces, social context, and adversarial or game-theoretic evaluation.

\section{Conclusion}

Representation Genesis treats better language as a first-class product of intelligent systems.
If WisdomBench measures whether agents learn from experience, Representation Genesis asks whether they can invent the representations that make future learning easier.

\begin{thebibliography}{9}
\bibitem{freedman2026compression}
Michael H. Freedman et al.
\newblock Compression is all you need: Modeling Mathematics.
\newblock arXiv:2603.20396, 2026.

\bibitem{ellis2020dreamcoder}
Kevin Ellis et al.
\newblock DreamCoder: Growing generalizable, interpretable knowledge with wake-sleep Bayesian program learning.
\newblock arXiv:2006.08381, 2020.

\bibitem{bowers2023stitch}
Matthew Bowers et al.
\newblock Top-Down Synthesis for Library Learning.
\newblock POPL, 2023.

\bibitem{grand2023lilo}
Gabriel Grand et al.
\newblock LILO: Learning Interpretable Libraries by Compressing and Documenting Code.
\newblock arXiv:2310.19791, 2023.

\bibitem{udrescu2020aifeynman}
Silviu-Marian Udrescu and Max Tegmark.
\newblock AI Feynman: A physics-inspired method for symbolic regression.
\newblock Science Advances, 2020.

\bibitem{schmidhuber2008compression}
Juergen Schmidhuber.
\newblock Driven by compression progress.
\newblock arXiv:0812.4360, 2008.
\end{thebibliography}

\end{document}
